\section{Arithmetic estimates for close pairs}
\label{sec:arithmetic}

Proposition~\ref{tab:conditional-reduction} leaves three comparison and
partition inputs to establish. We begin with the arithmetic estimates
needed for adjacent simultaneous births. These estimates concern pairs
of rational grids with fixed phases: residue averaging and a rotation
alternative control close pairs, while a weighted lcm bound describes
the denominator families that can concentrate when the determinant
interval is very short. The elementary determinant coordinates and
residue averages are in Lemma~\ref{ar:cycle}; all the estimates below
are finite and uniform in the phases.

\subsection{Small-prime exclusions}

Consider grids $(z+\alpha_1)/v$ and $(z'+\alpha_2)/w$, with positive integer
denominators $v,w$, fixed real phases, and positive half-widths $r_1,r_2$.
Write $g=(v,w)$ and $\eta=[v,w]\max(r_1,r_2)$. As in
Lemma~\ref{ar:cycle}, distance $O(\max(r_1,r_2))$ restricts the integer
determinant $(w/g)z-(v/g)z'$ to an interval of length $O(\eta)$.

We next record the Selberg upper-bound sieve in the form needed
for close grid pairs~\cite{Selberg1947,Uchiyama1962}. We give its
square-majorant and quadratic-form calculation in full, including
the determinant-cycle error bound. The same calculation will also
be used for prime values of a rotation parameter.

\begin{lemma}[Small-prime exclusions]\label{ar:sieve}
For every fixed $C>0$ there is an absolute sufficiently small $\xi>0$
with the following property. Suppose $\eta=[v,w]\max(r_1,r_2)\ge2$.
At any chosen primes $p\le\eta^\xi$, impose one forbidden numerator
residue on either or both grids, and impose a residue condition only
when $p$ divides that grid's denominator. The number of surviving pairs
at distance at most $C\max(r_1,r_2)$, per unit first-point translate, is
\[
 \ll_C g\eta\prod R_p,
\]
where the product has one factor for each imposed one-row exclusion.
\end{lemma}

\begin{proof}
Let $\mathcal P$ be the chosen primes, and let $g_0(p)$ be the probability
of a violation at $p$ after the complete residue average of
Lemma~\ref{ar:cycle}. It is $1/p$ for a one-row condition and
$2/p-1/p^2$ for a two-row condition. The violations are independent across
primes. Put
\[
 h_0'(p)=\frac{g_0(p)}{1-g_0(p)},\qquad
 Z=\eta^{1/100},\qquad
 \mathcal S=\sum_{\substack{d\le Z\text{ squarefree}\\p\mid d\Rightarrow p\in\mathcal P}}
 h_0'(d),
\]
extending $g_0,h_0'$ multiplicatively on squarefree integers.
All squarefree sums below use this support.

For coefficients $\lambda_1=1$ supported on $j\le Z$, the square
\[
 \left(\sum_j\lambda_j
 \mathbf1_{\{\text{a violation at every }p\mid j\}}\right)^2
\]
majorizes survival: at a surviving pair only $j=1$ contributes, giving
value one, and elsewhere the square is nonnegative. Its normalized mean is
\[
 \sum_{j,k}\lambda_j\lambda_k g_0([j,k])
 =\sum_d\frac1{h_0'(d)}
 \left(\sum_{d\mid j}\lambda_jg_0(j)\right)^2.
\]
Choose the inner sums to equal $\mu(d)h_0'(d)/\mathcal S$, where
$\mu(d)=(-1)^{\omega(d)}$ on the squarefree support. Divisor inversion
gives
\[
 \lambda_jg_0(j)=\frac{\mu(j)}{\mathcal S}
 \sum_{\substack{k\le Z/j\\(k,j)=1}}h_0'(jk).
\]
Thus $\lambda_1=1$, $|\lambda_j|\le1/g_0(j)\le j$, and the mean of the
majorant is $1/\mathcal S$. For a term indexed by $j,k$, the residue
moduli on both grids divide $[j,k]$. Lemma~\ref{ar:cycle}, followed by
the coefficient bounds, gives a total roundoff bounded by
\[
 Cg\sum_{j,k\le Z}|\lambda_j\lambda_k|
       [j,k]^2g_0([j,k])
 \le Cg\left(\sum_{j\le Z}j^3\right)^2
 \ll gZ^8\ll gZ^9.
\]
The inversion above takes place on a divisor-closed support, so its
$j=1$ numerator is exactly $\mathcal S$; every other absolute numerator
is at most $\mathcal S$. This also verifies the coefficient normalization
without extending the truncated support.

To estimate $\mathcal S$, normalize the full product
$\prod_{p\in\mathcal P}(1+h_0'(p))$ to a probability law on squarefree
products. Its expected logarithm is
\[
 \sum_{p\in\mathcal P}g_0(p)\log p
 \ll\xi\log\eta+1.
\]
For small enough $\xi$, Markov's inequality shows that a fixed positive
fraction of this product mass lies below $Z$. Consequently
\[
 \frac1{\mathcal S}\ll
 \prod_{p\in\mathcal P}(1-g_0(p)).
\]
The product on the right is the desired product of one-row factors. It
is bounded below by a fixed inverse power of $\log(2\eta)$, using the
elementary prime reciprocal bound. Hence $Z^9=\eta^{.09}$ is absorbed
by the main bound for large $\eta$. Bounded $\eta$ is absorbed in the
constant.
\end{proof}

\subsection{Rotation in an interval of integers or primes}

The next lemma turns an unexpectedly frequent close return into a
rational approximation. The target phase is arbitrary; no statistical
assumption on the phase is used.

\begin{lemma}[Rotation alternative]\label{ar:rotation}
There are absolute constants $C_1\gg C_2\gg1$ such that, for all
sufficiently large $u$, the following holds. Let $\alpha,y\in\mathbb R$,
$0\le\Delta\le e^{-C_1u}$ and $N\ge e^{C_1u}$. Either $\alpha$ lies
within $e^{C_2u}\Delta/N$ of a rational with reduced denominator at most
$e^{C_2u}$, or
\begin{align*}
 \#\{n\in[N,2N)\cap\mathbb Z:\|n\alpha-y\|\le\Delta\}
 &\ll e^{-u}N,\\
 \#\{p\in[N,2N)\text{ prime}:\|p\alpha-y\|\le\Delta\}
 &\ll e^{-u}N/\log N.
\end{align*}
\end{lemma}

\begin{proof}
Set $\varepsilon=e^{-10u}$. Pigeonholing gives
\[
 \alpha=\frac al+\theta,\qquad
 (a,l)=1,\qquad 1\le l\le\sqrt N,\qquad
 |\theta|\le\frac2{l\sqrt N}.
\]
We take $C_2$ of order $300$ and then $C_1$ sufficiently large.

\paragraph{Large denominators: $l>N^{1/30}$.}
For each integer $k\le N^{1/500}$, count multiples $n=km$ in the
interval, using blocks of $l$ consecutive values of $m$. The rational
parts $akm/l$ form an equally repeated grid of mesh $(k,l)/l\le k/l$.
Within a block the drift is $O(k/\sqrt N)$. Therefore the number with
$\|n\alpha-y\|\le\varepsilon$ is
\[
 \frac{2\varepsilon N}{k}+O(N^{.99}).
\]
Indeed complete-block errors total $O(N/l+\sqrt N)$, and an incomplete
block costs $O(l)$; all are within the displayed error.

For primes, use the square majorant from Lemma~\ref{ar:sieve} with the
bad condition $p\mid n$, primes $p\le N^\iota$, and
$Z=N^{1/5000}$, where $\iota>0$ is a sufficiently small absolute
constant. All tested divisors $[j,k]$ are within $N^{1/500}$. The full
product has mass at least a constant times $\log N$, by
$\prod_{p\le x}R_p^{-1}\ge\sum_{m\le x}1/m$; the same logarithmic
Markov argument retains a fixed fraction below $Z$. Thus the sieve
normalization is $\mathcal S\gg\log N$. The resulting prime count is
\[
 O\left(\frac{\varepsilon N}{\log N}
       +Z^4N^{.99}\right).
\]
Primes in $[N,2N)$ avoid every sieving prime. Both the displayed error
and the integer discrepancy are smaller than the required bounds when
$C_1$ is large.

\paragraph{Intermediate denominators:
$\varepsilon^{-10}<l\le N^{1/30}$.}
Divide $[N,2N)$ into windows of length comparable to $H=N^{1/4}$.
The drift in a window is $O(N^{-1/4})$. Inclusion-exclusion shows that
an interval of residue length $O(\varepsilon l)$ contains at most
\[
 O(\varepsilon\phi(l)+\tau(l))
 \ll\varepsilon^{1/2}\phi(l)
\]
invertible residues. For the last inequality, use the elementary
prime-wise bounds $\tau(l)\ll l^{1/8}$ and
$\phi(l)\gg l^{7/8}$, and then $l>\varepsilon^{-10}$.

Within one invertible progression modulo $l$ in a window, apply the
same square majorant, now omitting primes dividing $l$. Its main count
at a tested squarefree divisor $j$ is $H/(lj)+O(1)$. Removing the sieving
primes dividing $l$ multiplies the full product by a factor at least
$R(l)$, so
\[
 \frac1{\mathcal S}\ll\frac{R(l)^{-1}}{\log N}.
\]
The prime count in that progression is consequently
\[
 O\left(\frac{H}{\phi(l)\log N}+Z^4\right)
 \ll\frac{H}{\phi(l)\log N}.
\]
The last absorption is uniform for $l\le N^{1/30}$. Summing over
the possible invertible residues and the $O(N/H)$ windows gives
$O(\varepsilon^{1/2}N/\log N)$. For integers, direct residue counting
gives $O(\varepsilon^{1/2}N)$. Primes in the original interval are
coprime to $l$ because $l<N$.

\paragraph{Small denominators: $l\le\varepsilon^{-10}$.}
Suppose first that $|\theta|N>\Delta\varepsilon^{-20}$. For a fixed
residue $r\bmod l$, the continuous condition
\[
 \|ar/l+\theta x-y\|\le\Delta
\]
on $[N,2N)$ has $O(|\theta|N+1)$ components. Their total length is
$O(\Delta N+\Delta/|\theta|)$. Enlarging to the union of $H$-windows
meeting them gives total length
\[
 O\left((\Delta+|\theta|H)N+
              \frac{\Delta}{|\theta|}+H\right)
 \ll\varepsilon N.
\]
Here $|\theta|H\ll N^{-1/4}$,
$\Delta/|\theta|<N\varepsilon^{20}$, and $N$ is sufficiently large
relative to $\varepsilon^{-1}$. Each selected window has the
progression bound just proved. Summing first over windows and then over
the $\phi(l)$ invertible residues cancels that denominator and gives
$O(\varepsilon N/\log N)$ primes. Counting all $l$ residues instead
gives $O(\varepsilon N)$ integers.
In detail, if the selected windows for residue $r$ have total length
$L_r\ll\varepsilon N$, their number is $O(L_r/H)$. Integer counts per
window are at most $H/l+1\ll H/l$, because $H\gg l$. Prime counts per
window are $O(H/(\phi(l)\log N))$: the error $Z^4=N^{.0008}$ is
absorbed uniformly since $l\le N^{1/30}$ and
$H/(\phi(l)\log N)\ge N^{1/4-1/30}/\log N$. Thus each residue costs
$O(\varepsilon N/l)$ or $O(\varepsilon N/(\phi(l)\log N))$,
respectively, even though its selected windows depend on $r$.

If the drift inequality fails, then
\[
 l\le e^{100u},\qquad
 |\theta|\le\Delta e^{200u}/N,
\]
which is the rational alternative. These cases exhaust the possibilities,
and $\varepsilon^{1/2}=e^{-5u}$ is more than the stated saving.
When $\Delta=0$ and $\theta\ne0$, the drift components can be isolated
points; their $H$-window enlargement has the same bound. When
$\theta=0$, the small-denominator case is already the exact rational
alternative, and the other two denominator cases use the same rational
grid counts as above.
\end{proof}

\subsection{A weighted least-common-multiple estimate}

We now isolate the arithmetic structure of denominator pairs.
The prime-stripping induction and common-pivot reduction follow the
structural approach of Green and Walker~\cite[Proposition~1.2 and
Section~3]{GW2021} and Hauke-Treuer, Vazquez and
Walker~\cite[Proposition~2.1]{HVW2026}. The lcm cutoff, totient-bounded
vertex weights and moment-controlled edge inflation below require
the stated local argument. For related unweighted small-lcm counts
on dyadic integer sets, see Gou~\cite[Corollary~1.4 and
Theorem~1.5]{Gou2026}. Here we allow arbitrary finite supports and
weighted edges, with the following inflation assumptions. Let
$h_p\ge1$ be bounded, tend to one, and satisfy
\begin{equation}\label{ar:h-moments}
 \prod_p\left(1+\frac{h_p^m-1}{p}\right)<\infty
 \qquad(m>0).
\end{equation}
Put $h(n)=\prod_{p\mid n}h_p$. Uniformity in a family of such functions
means a uniform bound for $\sup_ph_p$ and for every fixed moment product
used below. Constants may depend on these bounds. In particular, for
each $\epsilon>0$, $h(n)\ll_\epsilon n^\epsilon$ uniformly in such a
family.

A weighted bipartite graph consists here of finite sets of positive
integers, vertex weights $f(v),g(w)$, and nonnegative edge weights
$e(v,w)$. The two copies of an integer are distinct vertices. Write
\[
 F=\sum_vf(v),\qquad G=\sum_wg(w),\qquad
 \mathcal E=\sum_{v,w}e(v,w).
\]

\begin{proposition}[Lcm bound]\label{ar:lcm}
Suppose $f(v)\le\phi(v)$, $g(w)\le\phi(w)$, and
\[
 e(v,w)\le f(v)g(w)
 h\bigl(v/(v,w)\bigr)h\bigl(w/(v,w)\bigr).
\]
If every positive edge satisfies $[v,w]\le Z$, then, for each fixed
$1/2<u<1$,
\begin{equation}\label{eq:lcm}
 \mathcal E\ll Z^{2(1-u)}(FG)^u.
\end{equation}
More generally, in a cell fixing the exponents $i_p,j_p$ at a specified
set of primes, its edge total satisfies the same bound with its own
vertex totals and the additional factor
\[
 h(K)K^{-(1-u)},\qquad K=\prod_p p^{|i_p-j_p|}.
\]
\end{proposition}

\begin{proof}
The cases $FG=0$ are immediate.

\paragraph{Fixed-exponent stripping.}
Fixing exponents at a prime and stripping
its powers divides the lcm cutoff by $p^{\max(i,j)}$. Divide the two
vertex weights by $\phi(p^i),\phi(p^j)$ and the edge weight by their
product and by $h_p^{\mathbf1_{i\ne j}}$. The stripped graph meets the
same hypotheses. Rescaling a proposed bound costs at most
\[
 h_p^{\mathbf1_{i\ne j}}
 p^{-2(1-u)\max(i,j)}
 \bigl(\phi(p^i)\phi(p^j)\bigr)^{1-u}
 \le h_p^{\mathbf1_{i\ne j}}p^{-(1-u)|i-j|}.
\]
Iteration proves the asserted cell coefficient once the main bound is
known.

It suffices first to prove the bound for graphs with no prime below a
fixed large cutoff $p_0$. Indeed, after that proof, strip the finitely
many smaller primes and sum their exponent cells. Each corresponding
kernel $h_p^{\mathbf1_{i\ne j}}p^{-(1-u)|i-j|}$ has a finite
convolution bound on the $u$-powers of probability vectors. Their
finite product only changes the constant.

\paragraph{Large-prime induction and common pivots.}
For the large-prime graph use induction on the number of occurring
primes, with a constant to be fixed below. The zero-prime graph is
immediate. If the assertion with this constant fails, normalize the
edge weights to a probability law. At each occurring prime the
induction hypothesis on a fixed exponent cell gives
\begin{equation}\label{eq:prime-cell}
 m_{ij}\le p^{-(1-u)|i-j|}
 h_p^{\mathbf1_{i\ne j}}(\alpha_i\beta_j)^u,
\end{equation}
where $\alpha_i,\beta_j$ are the fractions of the two vertex totals in
those exponent classes. Lemma~\ref{ar:pivot} gives a common exponent
$k_p$ with total off-pivot vertex mass $O(1/p)$. The edge probability
of two off-pivot exponents is $O(p^{-2u})$, and that of a single
deviation of size at least two is
\[
 O\bigl(p^{-2(1-u)-u}\bigr)=O(p^{-(2-u)}).
\]
Both sums converge. Increasing $p_0$ therefore permits us to discard
less than half the total edge weight while excluding both events at
every prime.

Put $N=\prod_pp^{k_p}$. Every remaining edge has the form
\[
 v=Na/b,\qquad w=Nc/d_1,\qquad [v,w]=Nac,
\]
where $b,d_1\mid N$ and $a,b,c,d_1$ are pairwise coprime squarefree
integers. These statements are edgewise: at a prime at most one side
differs from the pivot, and its difference is exactly one.
Since $(a,b)=(c,d_1)=1$, each vertex has a unique reduced ratio to $N$;
this remains true when $a$ or $c$ shares primes with $N$. Moreover,
\[
 (v,w)=N/(bd_1),\qquad
 h\bigl(v/(v,w)\bigr)h\bigl(w/(v,w)\bigr)=h(ab)h(cd_1).
\]

\paragraph{Dyadic capacities and summation.}
For $a\in[A,2A)$, set
\[
 C=\frac Z{NA},\qquad D_1=NA^2,\qquad D_2=NC^2.
\]
There are no such edges unless $C\ge1$. For each fixed $m>0$, the total
upper vertex weight on the first side, inflated by $h(ab)^m$, is
$O_m(D_1)$. To see this, use
\[
 \phi(Na/b)\le a\phi(N/b),\qquad
 \sum_{b\mid N}\frac{\phi(N/b)}N h(b)^m
 \le\prod_{p\mid N}\left(1+\frac{h_p^m-1}{p}\right),
\]
and expand $h(a)^m$ into its nonnegative squarefree divisor sum.
Divisibility by $j$ among $a\le2A$ occurs at most $2A/j$ times.
Equation~\eqref{ar:h-moments} proves the bound. The same calculation
on $c\le C$ gives $O_m(D_2)$.

Holder's inequality now bounds the retained edges in this dyad by
\[
 O_m\left((D_1D_2)^{1/m}
   \bigl(\min(F,D_1)\min(G,D_2)\bigr)^{1-1/m}\right).
\]
Indeed apply Holder under the actual product measure $f(v)g(w)$.
Its mass on the dyad is at most
$\min(F,O(D_1))\min(G,O(D_2))$, and the $m$th inflation moment under
that measure is $O_m(D_1D_2)$ by the totient capacities. These constants
depend only on the prescribed moment bounds, not on $N$ or its prime
support.
Choose $1-1/m>u$. Since $D_1D_2=Z^2$, divide this expression by
$Z^{2(1-u)}(FG)^u$ and put $f_1=F/D_1$, $g_1=G/D_2$. The quotient is
at most a constant times
\[
 \frac{\min(f_1,1)^{1-1/m}}{f_1^u}
 \frac{\min(g_1,1)^{1-1/m}}{g_1^u}
 \le\min(f_1,f_1^{-1})^{c(u)}
\]
for some $c(u)>0$. This is summable over dyadic $A$, since $D_1=NA^2$
changes geometrically. Thus the retained edges satisfy
\eqref{eq:lcm} with a constant independent of the number of primes.
Choose the inductive constant larger than twice this constant. Since
at least half the original edge weight was retained, this contradicts
the supposed failure and closes the induction. Stripping small primes
as above proves the unrestricted bound and its cell form.
\end{proof}

\subsection{Structure when the determinant scale is small}

For the application, vertex weights are multiplied by spatial widths.
Thus, from now on in this subsection, start with the graph in
Proposition~\ref{ar:lcm}, multiply the first and second vertex weights
by $r_1',r_2'>0$, and multiply the edges by $r_1'r_2'$. Reuse
$F,G,\mathcal E$ for these normalized totals, and put
\[
 r_+'=\max(r_1',r_2'),\qquad
 \mathcal R=\frac{\max(r_1',r_2')}{\min(r_1',r_2')}.
\]
At determinant scale $0<\eta<1$, meaning that the edge lcms are
comparable to $\eta/r_+'$, the ordinary bound supplies a factor
$\eta^{2(1-u)}$. Since $2(1-u)<1$, this alone does not give a summable
saving after a pair count costs $1/\eta$. We shall show that failure
of a stronger estimate forces a structured family of denominator pairs.
The following alternative is the structural output of the section.

\begin{proposition}[Small-scale extraction]\label{ar:extraction}
Fix positive constants $c_-,c_+$. Suppose $F,G\le1$ and every positive
edge satisfies
\[
 c_-\eta/r_+'\le[v,w]\le c_+\eta/r_+',\qquad 0<\eta<1.
\]
Given any sufficiently small fixed $\vartheta>0$, there are
$u\in(1/2,1)$, $\chi>0$, and $\eta_0>0$ such that for $\eta<\eta_0$
either
\begin{equation}\label{eq:small-saving}
 \mathcal E\le\eta^{1+\chi}\mathcal R^{-1/10}(FG)^u,
\end{equation}
or $\mathcal R\le\eta^{-\vartheta}$ and a subfamily of edges has the
form
\[
 v=NU_0a/b,\qquad w=NT_0c/d_1,\qquad
 [v,w]=NU_0T_0ac.
\]
Here $N,U_0,T_0$ are fixed positive integers,
\[
 (U_0,T_0)=1,\qquad U_0T_0\le\eta^{-\vartheta},
\]
$b,d_1\mid N$, and $a,b,c,d_1$ are pairwise coprime squarefree integers
coprime to $U_0T_0$. Moreover, $a\in[A,2A)$ and $c\in[C,2C)$ for fixed
dyadic $A,C$. The capacities
\[
 D_1=r_1'NU_0A^2,\qquad D_2=r_2'NT_0C^2
\]
satisfy $\eta^{2+\vartheta}\le D_i\le\eta^{-\vartheta}$, and the
captured normalized edge weight is at least $\eta^\vartheta D_1D_2$.
Every prime in the parameters occurs in the original lists.
\end{proposition}

\begin{proof}
We may suppose $\mathcal E>0$ and that \eqref{eq:small-saving} fails.
We shall first fix exceptional exponent coordinates, then concentrate
the remaining coordinates at common pivots, and finally select one
occupied pair of dyadic ranges. Small losses in powers of $\eta$ will
be chosen within a reserve much smaller than $\vartheta$.

\paragraph{Choosing a cell of maximal quality.}
Take $u=1/2+\epsilon_0<.52$, with $\epsilon_0>0$ small, and set
$s=1-u-.02$. For a cell fixing joint exponents at any subset of the
occurring primes, let $F',G',\mathcal E'$ be its normalized totals and
let $K$ be the product of its fixed exponent ratios. Call primes
$p\le p_0$ small. Among all cells of positive edge weight choose one
maximizing
\[
 \frac{\mathcal E'}{(F'G')^u}
 K^s H^{\#\{\text{small fixed primes}\}},
\]
where $H$ is a sufficiently large fixed number. A maximum exists since
the original supports are finite.

Every occurring small prime is fixed in this cell. Otherwise maximality
would bound the edge probability after further fixing it by
\[
 H^{-1}p^{-s|i-j|}(\alpha_i\beta_j)^u.
\]
Its sum is less than one when $H$ exceeds the bounded optimal constants
in the one-prime estimates
for $p\le p_0$, a contradiction. Here and below $\alpha,\beta$ denote
normalized vertex distributions inside the selected cell; they are not
assumed to be edge marginals.

In normalized variables, Proposition~\ref{ar:lcm} and its cell version
give
\[
 \mathcal E'\ll
 \eta^{2(1-u)}\mathcal R^{-(1-u)}
 h(K)K^{-(1-u)}(F'G')^u.
\]
Compare this with the quality inherited from the unrestricted graph and
the failure of \eqref{eq:small-saving}. It follows that
\begin{equation}\label{ar:fixed-ratio}
 K^{.02}h(K)^{-1}\mathcal R^{1-u-1/10}
 \ll_{p_0}\eta^{-2\epsilon_0-\chi}.
\end{equation}
Since $h(K)\ll_\epsilon K^\epsilon$, both $K$ and $\mathcal R$ are
bounded by arbitrarily small powers of $1/\eta$ on choosing
$\epsilon_0,\chi$ small and subsequently $\eta$ small.

Let $m$ be the normalized edge law in the selected cell. At a remaining
prime, maximality gives the bound of Lemma~\ref{ar:pivot} with price
$p^{-s|i-j|}$, even without an $h_p$ factor. On any joint subset of
remaining primes, the proved cell bound additionally gives
\begin{equation}\label{eq:joint-cell}
 m_{ij}\ll_{p_0}\eta^{-2\epsilon_0-\chi}
 K_{ij}^{-(1-u)}h(K_{ij})(\alpha_i\beta_j)^u.
\end{equation}
In this display $i,j$ are exponent vectors on the queried subset and
$K_{ij}$ is their ratio product. To obtain it, divide the combined-cell
upper bound by the quality lower bound for $\mathcal E'$; the fixed
factor $h(K)K^{-.02}$ and the favorable power of $\mathcal R$ are
bounded.

The weaker prime price $p^{-s|i-j|}$ does not give the summable
double-deviation bounds used in the proof of Proposition~\ref{ar:lcm}.
Instead of discarding all such deviations, we will fix one common
exceptional pattern. The next estimates control both its logarithmic
ratio, which bounds the factors $U_0T_0$, and the negative logarithm
of its probability, which bounds the mass lost when the pattern is fixed.

\paragraph{Controlling exceptional exponent patterns.}
For a remaining prime let $k$ be its pivot, let
$d=(1-\alpha_k)+(1-\beta_k)$, and set
\[
 x=m(i\ne k)+m(j\ne k).
\]
The central majorant implies $x\ge c_ud$, and Lemma~\ref{ar:pivot}
gives
\[
 x\ll p^{-s/(1-u)},\qquad
 B_p':=m(i\ne k,j\ne k)\ll x^{2u}.
\]
Explicitly,
\[
 x\ge1-m_{kk}\ge1-(\alpha_k\beta_k)^u\ge c_ud.
\]
Fix $1/2<u'<u$. Summing the $u'/u$ powers of the double-off majorant
gives, by convolution and Cauchy,
\[
 \sum_{i,j\ne k}m_{ij}^{u'/u}
 \ll\|\alpha_{\ne k}^{u'}\|_2\|\beta_{\ne k}^{u'}\|_2
 \le\bigl((1-\alpha_k)(1-\beta_k)\bigr)^{u'}
 \ll x^{2u'}.
\]
Since
$-t\log t\ll t^{u'/u}$ for $0\le t\le1$, its entropy is also
$O(x^{2u'})$. Inserting $|i-j|$ in the geometric kernel proves that
its contribution to $\mathbb E_m|i-j|$ is $O(x^{2u'})$.

Choose $\epsilon'>0$ so small that $(1-\epsilon')2u'>1$. For primes
with $x\le p^{-1+\epsilon'}$, all costs $x^{2u'}\log p$ are summable.
For the other primes we claim
\begin{align}
 \operatorname{Ent}_m(i)+\operatorname{Ent}_m(j)
 &\le x\log(1/x)+O(x)+O(x^{2u'}),\label{ar:marginal-entropy}\\
 \mathbb E_m|i-j|&\ge x-2B_p'.\label{ar:ratio-expectation}
\end{align}
Entropy means $-\sum_aP(a)\log P(a)$, with the summand zero at a zero
atom. For the first inequality, single deviations at distance at most
four occupy a bounded number of offsets, of total mass at most $x$.
Single deviations at distance at least five have entropy
$O(x^{2u'})$: indeed the sum of the $.8$ powers of their geometric
majorants is $O(p^{-4s})$, and $4s>2u'(1-\epsilon')$ for the choices
above. Double-off entropy was already bounded. Passing from entries
to marginal atoms only decreases these entropy upper bounds. The second
inequality follows because each single deviation contributes at least
one, while double deviations are counted twice in $x$.
For clarity, the leading entropy term is not counted twice. Set
$x_1=m(i\ne k)$ and $x_2=m(j\ne k)$. The bounded-offset part of the
two marginals has entropy at most
\[
 x_1\log(1/x_1)+x_2\log(1/x_2)+O(x)
 \le x\log(1/x)+O(x),\qquad x_1+x_2=x.
\]
The unbounded-offset contributions are the double-off and far-single
errors already estimated. In the far-single estimate, the explicit
geometric sum is
$O(\sum_{a\ge5}p^{-.8sa})=O(p^{-4s})$.

\paragraph{An entropy budget across primes.}
We now combine the coordinate bounds without assuming independence
among primes. Apply \eqref{eq:joint-cell} on the joint subset with
$x>p^{-1+\epsilon'}$, take logarithms, and average under $m$.
The cross-entropy against $\alpha$ is at least the entropy of the first
edge marginal, and similarly for $\beta$. The joint edge entropy is at
most the sum of its two vector marginal entropies, which in turn are
at most their coordinate entropy sums. Thus
\begin{align*}
 (1-u)\mathbb E_m\log K_{ij}
 &\le (1-u)\sum_p
   \bigl(\operatorname{Ent}_m(i_p)+\operatorname{Ent}_m(j_p)\bigr)\\
 &\quad+\mathbb E_m\log h(K_{ij})
       +O_{p_0}(1)+(2\epsilon_0+\chi)|\log\eta|.
\end{align*}
Indeed, if $\mu,\nu$ are the two vector edge marginals, the entropy
expression obtained by averaging the logarithm is
\[
 \operatorname{Ent}(m)-u\operatorname{CE}(\mu,\alpha)
                         -u\operatorname{CE}(\nu,\beta)
 \le(1-u)\bigl(\operatorname{Ent}(\mu)+\operatorname{Ent}(\nu)\bigr).
\]
Here $\operatorname{CE}(\mu,\alpha)=-\sum_i\mu_i\log\alpha_i$.
If a vertex-law atom is zero, the majorant forces the corresponding
edge marginal atom to be zero, so all cross-entropy terms are taken on
their positive support. No equality of vertex and edge laws, or
independence among prime coordinates, is used.
Here $\log(1/x)\le(1-\epsilon')\log p$. Equations
\eqref{ar:marginal-entropy}--\eqref{ar:ratio-expectation} therefore
leave a positive multiple of $\epsilon'\sum_px\log p$ on the left.
For sufficiently large $p_0$, the $h$ term, $B_p'\log p$, and the
$O(x)$ terms are absorbed into that margin. Consequently
\begin{equation}\label{ar:entropy-budget}
 \sum_{p:\,x>p^{-1+\epsilon'}}x\log p
 \ll_{u,\epsilon'}|\log\eta|.
\end{equation}
The proportional constant is independent of a further increase of
$p_0$, after decreasing $\eta_0$ to absorb its additive constant.
The choices occur in this order: first $\epsilon_0,\chi$, hence $u$;
then $u'\in(1/2,u)$; then $\epsilon'$ with $(1-\epsilon')2u'>1$;
then $p_0$; and finally $\eta_0$. The constant depending on
$u,\epsilon'$ may be large, but is fixed before $p_0$ is increased.

\paragraph{Fixing the exceptional pattern.}
The purpose of \eqref{ar:entropy-budget} is to control the entire
exceptional pattern, not only each prime separately. Record $(i_p,j_p)$
when both differ from the pivot, and record a blank otherwise. Its
coordinate entropy, including the blank atom, is $O(x^{2u'})$; its
expected log-ratio is $O(x^{2u'}\log p)$. On the small-$x$ primes these
costs have an arbitrarily small tail. On the remaining primes they are
an arbitrarily small multiple of the budget
\eqref{ar:entropy-budget}, by increasing $p_0$ and using
$x\ll p^{-s/(1-u)}$. Subadditivity therefore makes both the full pattern
entropy and its expected log-ratio as small a multiple of
$|\log\eta|$ as desired.

Single deviations of distance at least two have total probability
bounded by the tail of
\[
 \sum_p O\bigl(p^{-2s-us/(1-u)}\bigr).
\]
Its exponent is greater than one for $u$ sufficiently close to $1/2$.
Discard these deviations. Markov's inequality for the double-pattern
log-ratio and its information $-\log\Pr(\text{pattern})$ leaves a
fixed positive probability on patterns of log-ratio at most
$\rho|\log\eta|$ and information at most $\rho|\log\eta|$, where
$\rho>0$ can be as small as desired. There are at most $\eta^{-\rho}$
such patterns. One pattern therefore retains at least a fixed constant
times $\eta^\rho$ of the original edge probability after the discards.

Fix it. At its primes and the previously fixed primes put the minimum
of the two exponents into $N$, and put the two surplus parts into
$U_0,T_0$. At every other prime put the pivot into $N$. The remaining
single-step deviations give $a,b,c,d_1$. They are squarefree, pairwise
coprime, and supported away from $U_0T_0$; also $b,d_1\mid N$.
This proves the required algebraic form. Equations
\eqref{ar:fixed-ratio} and the pattern log-ratio bound make
$\mathcal R,U_0T_0$ at most arbitrarily small powers of $1/\eta$.

\paragraph{Selecting an occupied dyadic box.}
We make the remaining small-power bookkeeping explicit in meaning.
In this paragraph $o_*(1)$ denotes an exponent that can be bounded in
absolute value by any prescribed positive reserve: first choose
$\epsilon_0,\chi$ small, then $p_0$ large for the pattern losses, and
finally $\eta$ sufficiently small. It is not a limit with changing
block parameters.

The inherited quality gives
\[
 \mathcal E'\ge\eta^{1+o_*(1)}(F'G')^u.
\]
Proposition~\ref{ar:lcm}, now with exponent $.8$, also gives
$\mathcal E'\ll\eta^{.4}(F'G')^{.8}$. Thus
$F'G'\ge\eta^{2+o_*(1)}$, and the structured family just obtained
has weight at least $\eta^{2+o_*(1)}$.
More explicitly, replacing the first small exponent by $\rho>0$ gives
\[
 F'G'\gg\eta^{(.6+\rho)/(.8-u)},\qquad
 \frac{.6+\rho}{.8-u}
 =2+\frac{\rho+2\epsilon_0}{.3-\epsilon_0}.
\]
This makes the arbitrary proximity of the resulting exponent to two
transparent; the retained-pattern probability only adds another
prescribed small exponent.

Sort it into dyads for $a,c$. The lcm constraint gives
\[
 AC\asymp\frac{\eta}{r_+'NU_0T_0},\qquad
 D_1D_2\asymp\frac{\eta^2}{\mathcal R U_0T_0}.
\]
For each $A$ only boundedly many $C$ occur. The moment-capacity proof in
Proposition~\ref{ar:lcm} applies with $D_1,D_2$, with additional
inflation at most $h(U_0T_0)\le\eta^{-o_*(1)}$. For example, if
$D_1\ge\eta^{-\vartheta/4}$, Holder and $F,G\le1$ give, for fixed
large $m$,
\[
 O\left(\eta^{-o_*(1)}D_1^{1/m}D_2\right)
 =O\left(\eta^{2-o_*(1)}D_1^{-(1-1/m)}\right).
\]
Summing this geometric tail, and the analogous tail for $D_2$, costs
$O(\eta^{2+c\vartheta-o_*(1)})$ for an absolute $c>0$.
It is negligible compared with the retained mass.

In all remaining boxes $D_i\le\eta^{-\vartheta/4}$, while the product
relation implies $D_i\ge\eta^{2+\vartheta}$ when the reserves are
sufficiently small. There are only $O(1+|\log\eta|)$ such boxes because
$D_1$ varies geometrically and there are boundedly many $C$ per $A$.
One box has weight at least $\eta^{2+o_*(1)}$. Since
$D_1D_2\ll\eta^2$, its weight exceeds
$\eta^\vartheta D_1D_2$ for sufficiently small reserves and $\eta$.
The upper capacity bound and the asserted bounds on
$\mathcal R,U_0T_0$ are weaker than the bounds already obtained.
Every parameter used only exponents of primes occurring in the lists.
\end{proof}

\subsection{Auxiliary sampling in an extracted box}

In later applications a vertex also carries phase or mask variables.
The extraction concerns the aggregate graph weights, but the following
upper-bound interpretation permits those variables to be retained.
Its probabilities belong to an unrestricted comparison law, not to the
normalized law of the extracted edges. The corollary transfers small
comparison probability into small weighted edge mass.

\begin{corollary}[Negligible subsets of a saturated box]
\label{ar:box-sampling}
Use a box given by Proposition~\ref{ar:extraction}. Suppose that each
vertex weight is obtained by integrating a nonnegative density bounded
by its normalized totient weight against a probability law of auxiliary
variables; restrictions may reduce this density. Suppose also that the
edge majorant holds before integration against the independent product
of the two auxiliary laws.

Sample $a,c$ uniformly from their dyadic integer intervals, sample
$b,d_1$ independently with law $\pi_N$, and then sample the auxiliary
variables with their unrestricted laws. Invalid or nonintegral parameter
choices may be assigned zero weight. An event of sampling probability
$O(\eta^{10\vartheta})$ carries
$o(\eta^\vartheta D_1D_2)$ of the normalized edge weight. The same is
true of edges whose first vertex belongs to a set of total normalized
vertex weight $O(\eta^{10\vartheta}D_1)$.
The auxiliary law on a side may be conditional on that side's vertex
parameters. This statement is only an event-mass bound: an application
that varies an integer factor while holding an auxiliary phase fixed
must separately verify the required independence or a uniform
conditional sampling law.
\end{corollary}

\begin{proof}
On valid choices,
\[
 \phi(NU_0a/b)\le U_0a\phi(N/b),\qquad
 \phi(NT_0c/d_1)\le T_0c\phi(N/d_1).
\]
Thus the uninflated normalized vertex measures are bounded by
$O(D_1)$ and $O(D_2)$ times their respective sampling laws. Their
product has total mass $O(D_1D_2)$. The second moment of the edge
inflation is $O(h(U_0T_0)^2D_1D_2)$ by the capacity calculation in
Proposition~\ref{ar:lcm}. Cauchy's inequality therefore bounds the
weight on an event of probability $O(\eta^{10\vartheta})$ by
\[
 O\left(\eta^{5\vartheta}h(U_0T_0)D_1D_2\right)
 =O(\eta^{4\vartheta}D_1D_2),
\]
using $h(U_0T_0)\ll\eta^{-\vartheta}$. This is negligible compared
with the saturation threshold.

For the vertex restriction, each structured vertex has a unique
reduced defect ratio $a/b$ or $c/d_1$ relative to its fixed base.
Its uninflated product mass is
$O(\eta^{10\vartheta}D_1D_2)$ by the restricted first mass and the
unrestricted second capacity. The same second-moment argument applies.
These are upper bounds on the extracted edges; they do not assert
independence of the actual surviving edge law.
\end{proof}

The rotation alternative controls close returns unless a phase is
well approximated by a rational. The extraction and sampling results
supply a family on which such returns can be tested by varying integer
factors. Section~\ref{sec:direct} applies these estimates to the birth
comparisons in Proposition~\ref{tab:first-join-input}.
